\begin{ouverture}[college2]{10}{Interactions et forces}{Partie 2 · Mouvements et interactions}{Pourquoi un objet posé ne tombe-t-il pas ?}
\objectif{Identifier les objets en interaction avec un système.}
\objectif{Distinguer contact et action à distance.}
\objectif{Représenter une force par quatre caractéristiques.}
\prerequis{Décrire le mouvement et choisir un système.}
\end{ouverture}
\pcPage{Le cours}{Interactions et forces}
\phantomsection\label{ess-10}
\essSection{1}{Délimiter le système}
Le système est l’objet ou l’ensemble d’objets que l’on choisit d’étudier. On recherche les objets extérieurs qui agissent sur lui. Un livre posé sur une table interagit notamment avec la Terre et la table. L’interaction avec la Terre est à distance ; celle avec la table est de contact. On indique si l’action de l’air est négligée dans le modèle.\par
\pcFigure{figures/c10-doi.pdf}{Diagramme objet-interactions du livre : trait plein pour un contact, pointillé pour une interaction à distance.}
\essSection{2}{Modéliser une action}
Une force représente l’action d’un objet sur un autre. Elle possède un point d’application, une direction, un sens et une valeur en newtons (N). La direction est la droite support ; le sens précise vers quel côté l’action s’exerce. Une flèche représente ces caractéristiques : origine, orientation, pointe et longueur donnée par une échelle. Le nom de l’auteur de l’action doit être identifiable.\par
\clearpage
\essSection{3}{Mesurer et dessiner}
Le dynamomètre mesure une valeur de force. Vérifier son zéro, son calibre et son alignement avec l’action. Pour dessiner une force de $4\unite N$ à l’échelle $1\unite{cm}$ pour $2\unite N$, tracer une flèche de $2\unite{cm}$. Une flèche longue ne donne pas une valeur si aucune échelle n’est annoncée.\par
\pcFigure{figures/c10-forces.pdf}{Pour le livre immobile, deux actions verticales opposées peuvent se compenser dans le modèle où les autres actions sont négligées.}
\essSection{4}{Ce que produit une action}
Une action peut déformer un objet ou modifier son mouvement : valeur de vitesse, direction ou sens. Un objet immobile peut subir plusieurs forces qui se compensent. Une force n’est donc pas nécessairement synonyme de mouvement visible. Les deux actions d’une interaction s’exercent sur deux objets différents : on ne les additionne pas comme des forces appliquées au même système.\par
\begin{pcEncadre}{method}{Un exemple expliqué}Un ballon est suspendu par un fil et immobile. Sur le ballon, on représente l’action de la Terre vers le bas et celle du fil vers le haut. Si leurs valeurs sont toutes deux $2\unite N$, à l’échelle $1\unite{cm}$ pour $1\unite N$, chaque flèche mesure $2\unite{cm}$. Ces forces sont appliquées au même ballon et se compensent dans le modèle.\end{pcEncadre}
\pcPage{Expérimenter et retenir}{Interactions et forces}
\phantomsection\label{activite-10}
\begin{pcEncadre}{method}{Investigation en classe}\textbf{Question.} Quelles actions s’exercent sur un objet suspendu ? Suspendre une masse scolaire à un dynamomètre adapté, sans dépasser le calibre. Relever la valeur lorsque l’objet est immobile. Construire le diagramme objet-interactions, puis représenter les forces avec une échelle choisie. Incliner légèrement le support uniquement selon les consignes : comparer la direction de l’instrument à celle de l’action mesurée. \textbf{À rendre :} système, interactions, mesure, flèches légendées et hypothèses.\end{pcEncadre}
\pcSousTitre{Ce que je dois savoir expliquer}\begin{questions}
\item Le bilan des actions dépend du système choisi.
\item Un contact et une interaction à distance sont deux catégories différentes.
\item Une force se représente par un point, une direction, un sens et une valeur.
\item Le newton est une unité de force ; le kilogramme est une unité de masse.
\item Un objet immobile peut subir plusieurs forces non nulles.
\end{questions}
\begin{pcEncadre}{attention}{Une erreur à éviter}La direction « verticale » et le sens « vers le bas » sont deux informations distinctes ; l’une ne remplace pas l’autre.\end{pcEncadre}
\pcRepere{cor-activite-10}{Pistes pour l’investigation}\pcGrille
\clearpage\markright{Je démarre}
\rubrique{Je démarre}
\par\noindent\begin{minipage}{\linewidth}
\begin{exercice}[1]{Choisir le système}
\begin{questions}
\item On étudie un ballon tenu par une main. Quel est le système ?
\item La main est-elle alors extérieure au système ?
\end{questions}
\end{exercice}
\end{minipage}
\par\noindent\begin{minipage}{\linewidth}
\begin{exercice}[1]{Contact ou distance}
\begin{questions}
\item Classer l’action d’une table sur un livre et celle de la Terre sur ce livre.
\item Un aimant doit-il toujours toucher un trombone pour agir ?
\end{questions}
\end{exercice}
\end{minipage}
\par\noindent\begin{minipage}{\linewidth}
\begin{exercice}[1]{Quatre caractéristiques}
\begin{questions}
\item Nommer les quatre caractéristiques d’une force.
\item Donner son unité.
\end{questions}
\end{exercice}
\end{minipage}
\par\noindent\begin{minipage}{\linewidth}
\begin{exercice}[1]{L’appareil adapté}
\begin{questions}
\item Quel appareil mesure la valeur d’une force ?
\item Pourquoi vérifier le calibre ?
\end{questions}
\end{exercice}
\end{minipage}
\par\noindent\begin{minipage}{\linewidth}
\begin{exercice}[1]{Deux sens}
\begin{questions}
\item Deux forces sont verticales, l’une vers le haut et l’autre vers le bas. Ont-elles la même direction ?
\item Ont-elles le même sens ?
\end{questions}
\end{exercice}
\end{minipage}
\clearpage\markright{Je m’entraîne}
\rubrique{Je m’entraîne}
\par\noindent\begin{minipage}{\linewidth}
\begin{exercice}[2]{Dessiner à l’échelle}
\begin{questions}
\item À $1\unite{cm}$ pour $2\unite N$, quelle longueur pour $6\unite N$ ?
\item Pour $1\unite N$ ?
\end{questions}
\end{exercice}
\end{minipage}
\par\noindent\begin{minipage}{\linewidth}
\begin{exercice}[2]{Lire une longueur}
\pcFigureExact{figures/ex-force-4cm.pdf}{Flèche de 4 cm à l’impression à 100 \% ; échelle : 1 cm pour 0,5 N.}
\begin{questions}
\item Une flèche mesure $4\unite{cm}$ à l’échelle $1\unite{cm}$ pour $0{,}5\unite N$. Quelle force représente-t-elle ?
\item Changer l’échelle change-t-il l’action réelle ?
\end{questions}
\end{exercice}
\end{minipage}
\par\noindent\begin{minipage}{\linewidth}
\begin{exercice}[2]{Un objet immobile}
\begin{questions}
\item Un objet immobile subit-il forcément une force nulle de la Terre ?
\item Comment peut-il rester immobile ?
\end{questions}
\end{exercice}
\end{minipage}
\par\noindent\begin{minipage}{\linewidth}
\begin{exercice}[2]{Le livre sur la table}
\begin{questions}
\item Construire le diagramme des interactions du livre, en négligeant l’air.
\item Quelles forces représenter sur le livre ?
\end{questions}
\end{exercice}
\end{minipage}
\par\noindent\begin{minipage}{\linewidth}
\begin{exercice}[2]{Une balle qui tourne}
\begin{questions}
\item Une balle peut-elle subir une action qui change sa direction sans augmenter la valeur de sa vitesse ?
\item Donner un exemple de mouvement compatible.
\end{questions}
\end{exercice}
\end{minipage}
\clearpage\markright{Je raisonne}
\rubrique{Je raisonne}
\par\noindent\begin{minipage}{\linewidth}
\begin{exercice}[3]{Deux objets, deux actions}
\begin{questions}
\item Une main pousse une boîte. La boîte agit-elle aussi sur la main ?
\item Ces deux actions se compensent-elles dans le bilan de la boîte ?
\end{questions}
\end{exercice}
\end{minipage}
\par\noindent\begin{minipage}{\linewidth}
\begin{exercice}[3]{Un bilan incomplet}
\pcFigureEx{figures/ex-suspension-incomplete.pdf}{Bilan de forces proposé par l’élève, incomplet.}
\begin{questions}
\item Pour un objet suspendu immobile, un élève ne dessine que la force du fil. Quelle action manque ?
\item Quelle erreur de conclusion pourrait-il faire ?
\end{questions}
\end{exercice}
\end{minipage}
\par\noindent\begin{minipage}{\linewidth}
\begin{exercice}[3]{Une lecture discutable}
\begin{questions}
\item Un dynamomètre affiche 0,2 N avant tout accrochage puis 2,2 N avec l’objet. Si le décalage reste constant, quelle correction proposer ?
\item Quelle précaution est préférable ?
\end{questions}
\end{exercice}
\end{minipage}
\par\noindent\begin{minipage}{\linewidth}
\begin{exercice}[3]{Une flèche sans légende}
\begin{questions}
\item Deux élèves tracent des flèches de longueurs différentes pour la même force. Lequel a tort ?
\item Quelles informations vérifier avant de décider ?
\end{questions}
\end{exercice}
\end{minipage}
\par\noindent\begin{minipage}{\linewidth}
\begin{exercice}[3]{Changer le système}
\begin{questions}
\item On étudie d’abord un objet suspendu, puis l’ensemble objet-fil. L’action du fil sur l’objet a-t-elle encore le même statut ?
\item Quelle action extérieure faut-il alors chercher ?
\end{questions}
\end{exercice}
\end{minipage}
\pcPage{Défi du chapitre}{Interactions et forces}
\phantomsection\label{defi-ex-10}
\pcDefiTitre{10}{Une suspension à expliquer}
\pcFigureEx{figures/ex-objet-suspendu.pdf}{}
\begin{questions}
\item Un objet est immobile sous un fil vertical. Le dynamomètre indique $3\unite N$. Définir le système et deux interactions.
\item Tracer les forces à l’échelle $1\unite{cm}$ pour $1{,}5\unite N$.
\item Un élève dit « aucune force puisque rien ne bouge ». Répondre.
\end{questions}
\pcGrille
